% ============================================================================
% Fichier  : partie4/P04_C09_vulnerabilite.tex
% Titre    : Qu'est-ce qu'une Vulnérabilité
% Auteur   : MINKA MI NGUIDJOÏ Thierry Emmanuel
% Version  : 1.0 -- novembre 2026
% Corrigés : corriges/Corr_C09.tex
% ============================================================================
% Sources : programme CVE (MITRE), CWE, NVD (NIST) ; CVSS v4.0 (FIRST, 2023)
% et v3.1 (2019) ; EPSS (FIRST) ; catalogue KEV (CISA) ; SSVC (CISA, 2020).
% Les scores CVSS v4.0 donnés (9,3 et 10,0) correspondent aux vecteurs cités
% (calculateur du FIRST) ; les autres valeurs sont fictives et marquées
% comme telles. Identifiants VULN-2025-xxxx fictifs.
% Faits nouveaux (bible) : vecteur de VULN-2024-4471 ; quatre
% vulnérabilités d'exemple chez SOLITRANS (V-A à V-D).
% ============================================================================

\chapter{Qu'est-ce qu'une Vulnérabilité}
\label{chap:vulnerabilite}
\chaptermark{Qu'est-ce qu'une vulnérabilité}

\epigraph{%
    « Un score dit la gravité d'une faille dans l'absolu.\\
    Il ne dit pas qu'elle sera exploitée chez vous. »%
}{Principe du cours}

\niveauChapitre{Fondamental}{%
    Lecture du Chapitre~\ref{chap:prologue}
}

\begin{boxobjectifs}
\begin{itemize}
    \item Définir une vulnérabilité et la distinguer de la faiblesse, de la
          menace et du risque.
    \item Lire un identifiant CVE et le rattacher à une faiblesse CWE.
    \item Lire un vecteur CVSS v4.0 et en expliquer chaque métrique.
    \item Interpréter un score EPSS et une inscription au catalogue KEV.
    \item Expliquer pourquoi la gravité technique ne suffit pas à décider
          d'une priorité.
\end{itemize}
\end{boxobjectifs}

\bigskip

% ============================================================================
\section*{Avant de commencer : 412 critiques}
% ============================================================================

Douala, mai 2025. SOLITRANS LOGISTIQUE\footnote{SOLITRANS LOGISTIQUE, ses
salariés, ses sites et ses chiffres sont fictifs~; les référentiels cités
(CVE, CWE, CVSS, EPSS, KEV, SSVC) sont exacts.}, transitaire et
logisticien du port, vient de réaliser son premier balayage complet de
vulnérabilités sur ses 1\,380 actifs. Le rapport de l'outil compte 3\,812
constats, dont 412 qualifiés de « critiques » par leur score. Mme~EBODE,
responsable de la gestion des vulnérabilités depuis trois semaines, doit
dire au comité de direction ce qu'on corrige d'abord. Le RSSI, M.~NJIKE, a
une réponse simple~: « les 412 critiques, dans l'ordre du score ». Mme~EBODE
sait que l'équipe système corrige environ quarante failles par semaine. Il
faudrait dix semaines pour les seules critiques.

\begin{boxalert}
\textbf{Violation typique.} Confondre la gravité technique d'une
vulnérabilité, mesurée par un score, avec le risque qu'elle fait courir à
l'organisation.

\smallskip
\textbf{Règle~:} le score de base décrit la faille telle qu'elle est partout.
La priorité de correction dépend aussi de ce que l'on sait de son
exploitation et de la place de l'actif dans l'organisation.
\end{boxalert}

\separateur

% ============================================================================
\section{Faiblesse, Vulnérabilité, Menace, Risque}
\label{sec:c09-definitions}
% ============================================================================

\begin{boxdef}
\textbf{\termecle{Faiblesse}}~: type d'erreur de conception ou de
programmation susceptible de créer des vulnérabilités, par exemple une
authentification mal vérifiée ou une entrée non contrôlée. Le catalogue CWE
(\textit{Common Weakness Enumeration}) les recense~\cite{mitrecwe}.

\smallskip
\textbf{\termecle{Vulnérabilité}}~: défaut présent dans une version précise
d'un produit, qui peut être exploité pour porter atteinte à la
confidentialité, à l'intégrité ou à la disponibilité. Le programme CVE
(\textit{Common Vulnerabilities and Exposures}) lui attribue un identifiant
unique~\cite{mitrecve}.

\smallskip
\textbf{\termecle{Menace}}~: source potentielle d'atteinte, par exemple un
groupe qui exploite une famille de failles.

\smallskip
\textbf{\termecle{Risque}}~: combinaison de la vraisemblance qu'une menace
exploite une vulnérabilité et de la gravité de l'atteinte qui en résulte
pour l'organisation.
\end{boxdef}

La faille de la passerelle de NYANGA CRÉDIT illustre la chaîne. La
faiblesse est un contrôle d'authentification incomplet (CWE-287,
authentification incorrecte). La vulnérabilité est son instance dans le
micrologiciel 7.2.1 de la passerelle, désignée dans ce livre par
l'identifiant fictif VULN-2024-4471. La menace, ce sont les groupes qui
l'exploitaient dès décembre 2024. Le risque, pour NYANGA, combinait une
passerelle exposée sur Internet, non corrigée, et un accès direct au
contrôleur de domaine.

\subsection{Le cycle de vie d'un identifiant CVE}

Un chercheur, un éditeur ou un utilisateur découvre une faille. Une
\termecle{autorité de numérotation} (CNA), souvent l'éditeur lui-même,
réserve un identifiant de la forme \texttt{CVE-AAAA-NNNNN}, puis le publie
avec une description quand le correctif ou l'avis est prêt. La base
nationale américaine NVD~\cite{nistnvd} enrichit ensuite l'enregistrement~:
produits et versions concernés, score, faiblesse CWE. Entre la publication
et l'enrichissement, plusieurs jours peuvent passer~; l'analyste qui attend
le score de la NVD perd ce temps.

% ============================================================================
\section{Le Score CVSS}
\label{sec:c09-cvss}
% ============================================================================

Le \termecle{CVSS} (\textit{Common Vulnerability Scoring System}), tenu par
le FIRST, mesure la gravité d'une vulnérabilité sur une échelle de 0 à 10.
Sa version~4.0, publiée en 2023~\cite{first2023cvss4}, remplace
progressivement la version~3.1~\cite{first2019cvss31}.

\subsection{Quatre groupes de métriques}

\begin{tableaucours}{L{2.8cm}L{5.4cm}L{4.0cm}}{%
    \tableauHeader{Groupe} &
    \tableauHeader{Ce qu'il mesure} &
    \tableauHeader{Qui le renseigne}
}
Base & Les propriétés de la faille, identiques partout & L'éditeur, la
NVD \\
Menace (\textit{Threat}) & La maturité de l'exploitation connue & L'analyste,
à partir du renseignement \\
Environnement & L'importance de l'actif et les mesures en place & L'organisation \\
Supplémentaire & Des informations sans effet sur le score (automatisable,
effort de reprise) & L'éditeur, l'analyste \\
\end{tableaucours}
\vspace{-0.8em}
\begin{center}{\small\itshape Tableau~9.1. Les groupes de métriques de CVSS v4.0}\end{center}

La version~4.0 exige que l'on dise quels groupes ont servi~: CVSS-B pour la
seule base, CVSS-BT avec la menace, CVSS-BE avec l'environnement, CVSS-BTE
avec les trois. Un « 9,3 » sans cette mention est ambigu.

\subsection{Les métriques de base}

\begin{tableaucours}{L{1.6cm}L{4.4cm}L{6.2cm}}{%
    \tableauHeader{Code} &
    \tableauHeader{Métrique} &
    \tableauHeader{Valeurs}
}
AV & Vecteur d'attaque & Réseau (N), adjacent (A), local (L), physique (P) \\
AC & Complexité de l'attaque & Faible (L), élevée (H) \\
AT & Conditions préalables (nouvelle en 4.0) & Aucune (N), présentes (P) \\
PR & Privilèges requis & Aucun (N), faibles (L), élevés (H) \\
UI & Interaction de l'utilisateur & Aucune (N), passive (P), active (A) \\
VC, VI, VA & Impact sur le système vulnérable (confidentialité, intégrité,
disponibilité) & Élevé (H), faible (L), nul (N) \\
SC, SI, SA & Impact sur les systèmes suivants & Élevé (H), faible (L), nul
(N) \\
\end{tableaucours}
\vspace{-0.8em}
\begin{center}{\small\itshape Tableau~9.2. Les métriques de base de CVSS v4.0}\end{center}

La distinction entre système vulnérable et systèmes suivants remplace la
notion de « portée » de la version~3.1. Elle permet de dire qu'une faille de
la passerelle n'atteint pas seulement la passerelle, mais aussi les systèmes
auxquels elle donne accès.

\begin{boxobjet}{Vecteur CVSS v4.0 de VULN-2024-4471}
\textbf{Vecteur publié par l'éditeur (fictif, construit pour le cours)~:}
\texttt{CVSS:4.0/AV:N/AC:L/AT:N/PR:N/UI:N/VC:H/VI:H/VA:H/SC:N/SI:N/SA:N}.

\textbf{Score CVSS-B~:} 9,3, critique.

\textbf{Lecture~:} exploitable depuis Internet, sans condition, sans
compte, sans action d'un utilisateur, avec un impact élevé sur la
passerelle elle-même.

\textbf{Ce que l'éditeur ne pouvait pas savoir~:} qu'à NYANGA, la passerelle
ouvrait un accès direct au contrôleur de domaine. L'analyste de NYANGA, en
renseignant SC, SI et SA à « élevé » dans le groupe environnemental, aurait
obtenu un score de 10,0 pour son propre contexte.
\end{boxobjet}

\subsection{Les seuils de gravité}

\begin{tableaucours}{L{3.4cm}L{3.4cm}}{%
    \tableauHeader{Score} &
    \tableauHeader{Gravité}
}
0,0 & Aucune \\
0,1 à 3,9 & Faible \\
4,0 à 6,9 & Moyenne \\
7,0 à 8,9 & Élevée \\
9,0 à 10,0 & Critique \\
\end{tableaucours}
\vspace{-0.8em}
\begin{center}{\small\itshape Tableau~9.3. Les seuils de gravité du CVSS}\end{center}

\begin{boxinfo}
\textbf{Un score ne se calcule pas de tête.} En version~4.0, le score se
déduit du vecteur par une table de correspondance publiée par le FIRST, et
non par une formule simple comme en version~3.1. On utilise le calculateur
officiel, et l'on note toujours le vecteur à côté du score~: c'est le
vecteur qui permet de discuter le score.
\end{boxinfo}

% ============================================================================
\section{EPSS : la Probabilité d'Exploitation}
\label{sec:c09-epss}
% ============================================================================

L'\termecle{EPSS} (\textit{Exploit Prediction Scoring System}), lui aussi
tenu par le FIRST~\cite{firstepss}, estime chaque jour, pour chaque CVE
publiée, la probabilité qu'elle fasse l'objet de tentatives d'exploitation
observées dans les trente jours suivants. Il s'appuie sur un modèle
statistique alimenté par des données d'exploitation réelles. Il donne une
probabilité, entre 0 et 1, et un centile, qui situe la CVE parmi toutes les
autres.

Le CVSS et l'EPSS répondent à deux questions différentes. Le premier dit
« si cette faille est exploitée, à quel point est-ce grave~? »~; le second
dit « quelle est la chance qu'elle soit exploitée bientôt~? ». Les données
publiées par le projet EPSS~\cite{firstepss} montrent qu'une faible part
des CVE publiées fait l'objet d'une exploitation observée~: un score
critique n'annonce donc pas une exploitation, et un score moyen ne
l'exclut pas.

% ============================================================================
\section{KEV : ce qui est Exploité}
\label{sec:c09-kev}
% ============================================================================

Le catalogue KEV (\textit{Known Exploited Vulnerabilities}), tenu par
l'agence américaine CISA~\cite{cisakev}, recense les vulnérabilités dont
l'exploitation est constatée. Une inscription n'est pas une prévision~:
c'est un fait. Le catalogue fixe une échéance de correction pour les
agences fédérales américaines, mais toute organisation peut s'en servir
comme d'une liste de failles à traiter sans attendre.

\begin{tableaucours}{L{2.4cm}L{4.2cm}L{3.0cm}L{2.6cm}}{%
    \tableauHeader{Indicateur} &
    \tableauHeader{Question} &
    \tableauHeader{Nature} &
    \tableauHeader{Fréquence}
}
CVSS-B & Quelle gravité si exploitée~? & Mesure technique & À la
publication \\
EPSS & Quelle chance d'exploitation dans 30 jours~? & Probabilité estimée
& Quotidienne \\
KEV & Est-elle déjà exploitée~? & Fait constaté & À l'inscription \\
\end{tableaucours}
\vspace{-0.8em}
\begin{center}{\small\itshape Tableau~9.4. Trois indicateurs, trois questions}\end{center}

% ============================================================================
\section{SSVC : Décider plutôt que Noter}
\label{sec:c09-ssvc}
% ============================================================================

La CISA propose aussi une méthode de décision, le SSVC
(\textit{Stakeholder-Specific Vulnerability Categorization})~\cite{cisa2022ssvc}.
Au lieu d'un score, elle conduit, par un arbre de questions, à l'une de
quatre décisions~: suivre (\textit{Track}), suivre de près
(\textit{Track*}), traiter rapidement (\textit{Attend}), agir sans délai
(\textit{Act}). Les questions portent sur l'état de l'exploitation (aucune,
preuve de concept, active), la possibilité d'automatiser l'attaque,
l'impact technique, et l'importance de l'actif pour la mission.

Le SSVC rend explicite ce que le score cache~: la décision dépend de
l'organisation. La même faille peut appeler « agir » sur une passerelle
exposée et « suivre » sur un serveur de test isolé.

% ============================================================================
\section{Quatre Failles chez SOLITRANS}
\label{sec:c09-solitrans}
% ============================================================================

\begin{boxscenario}{Opération SOLITRANS. Quatre lignes du rapport de balayage}
\begin{tableaucours}{L{1.0cm}L{3.8cm}cccL{1.8cm}}{%
    \tableauHeader{Réf.} &
    \tableauHeader{Actif} &
    \tableauHeader{CVSS-B} &
    \tableauHeader{EPSS} &
    \tableauHeader{KEV} &
    \tableauHeader{Exposé}
}
V-A & Serveur d'impression, réseau interne & 9,8 & 0,002 & Non & Non \\
V-B & Portail des transitaires, Internet & 7,5 & 0,640 & Oui & Oui \\
V-C & Base de suivi des conteneurs, interne & 8,8 & 0,031 & Non & Non \\
V-D & Pare-feu du site de Kribi, Internet & 6,5 & 0,410 & Oui & Oui \\
\end{tableaucours}

\smallskip
{\footnotesize Valeurs fictives, construites pour l'exercice. EPSS~:
probabilité estimée d'exploitation dans les trente jours.}

\medskip
\textbf{Constatations.} Rangées par score de base, les failles donnent
l'ordre V-A, V-C, V-B, V-D. Rangées par exploitation constatée puis par
probabilité, elles donnent V-B, V-D, V-C, V-A.

\textbf{Analyse} (confiance élevée). Les deux failles exploitées, sur des
actifs exposés à Internet, passent avant les deux autres, malgré leurs
scores inférieurs. V-A, la plus « grave » sur le papier, touche un serveur
d'impression interne et n'est pas exploitée~: elle reste à corriger, mais
pas cette semaine.
\end{boxscenario}

\begin{boxlecture}{Que sait-on de chaque faille~?}
\begin{tableaucours}{L{5.6cm}cccc}{%
    \tableauHeader{Question} &
    \tableauHeader{V-A} & \tableauHeader{V-B} & \tableauHeader{V-C} &
    \tableauHeader{V-D}
}
La faille est exploitée dans le monde & \NEK{} & \OUI{} & \NEK{} & \OUI{} \\
L'actif est joignable depuis Internet & \NON{} & \OUI{} & \NON{} & \OUI{} \\
La faille est présente (constat vérifié) & \NEK{} & \OUI{} & \OUI{} &
\OUI{} \\
\end{tableaucours}

\medskip
\textbf{Lecture.} Une absence d'inscription au catalogue KEV n'est pas une
preuve de non-exploitation~: V-A et V-C sont \NEK{}, pas \NON{}. Et V-A n'a
pas encore été vérifiée~: l'outil l'a déduite de la version annoncée par le
service, ce qui peut produire un faux positif (Chapitre~\ref{chap:identifier}).
\end{boxlecture}

% ============================================================================
\section{Ce que la Notion demande à chaque Métier}
\label{sec:c09-metiers}
% ============================================================================

\begin{boxmetier}{Analyste en gestion des vulnérabilités}
\textbf{Lire une vulnérabilité}
\begin{enumerate}
    \item Identifier la CVE, la faiblesse CWE et les versions touchées.
    \item Lire le vecteur CVSS, pas seulement le score, et noter la version
          du CVSS.
    \item Relever l'EPSS du jour et l'inscription éventuelle au catalogue
          KEV.
    \item Situer l'actif~: exposition, importance, mesures en place.
    \item En déduire une décision, de type SSVC, et la justifier.
\end{enumerate}
\end{boxmetier}

\begin{boxmetier}{RSSI}
Fixer la règle de priorité par écrit, en combinant exploitation, exposition
et importance de l'actif, et la faire approuver. Une règle fondée sur le
seul score conduit l'équipe à corriger le mauvais serveur en premier.
\end{boxmetier}

\begin{boxmetier}{Analyste en réponse à incident}
Pendant une enquête, le catalogue KEV et l'EPSS orientent la recherche du
point d'entrée~: une faille exploitée, présente sur un actif exposé, est la
première hypothèse à tester.
\end{boxmetier}

\separateur

\begin{boxresume}
\textbf{Ce qu'il faut retenir de ce chapitre~:}
\begin{itemize}
    \item La faiblesse (CWE) est un type d'erreur~; la vulnérabilité (CVE)
          son instance dans un produit~; le risque dépend de la menace et
          du contexte.
    \item Le CVSS v4.0 a quatre groupes de métriques~; on précise CVSS-B,
          -BT, -BE ou -BTE, et l'on donne toujours le vecteur.
    \item L'EPSS estime la probabilité d'exploitation dans les trente
          jours~; le catalogue KEV constate l'exploitation.
    \item Le SSVC remplace le score par une décision qui dépend de
          l'organisation.
    \item Chez SOLITRANS, l'ordre par score et l'ordre par exploitation et
          exposition s'inversent.
\end{itemize}
\end{boxresume}

\begin{boxinfo}
\textbf{Points de contrôle.} Vous avez compris ce chapitre si vous savez~:
\begin{itemize}
    \item expliquer chaque élément d'un vecteur CVSS v4.0~;
    \item dire ce que l'EPSS et le catalogue KEV ajoutent au score~;
    \item justifier un ordre de correction différent de l'ordre des scores.
\end{itemize}
\end{boxinfo}

% ============================================================================
\section*{Exercices}
% ============================================================================

\begin{boxexo}[title={\ding{45}\quad Exercice~9.1. Faiblesse ou vulnérabilité \niveaubase}]
Classez chaque élément comme faiblesse, vulnérabilité, menace ou risque~:
(a) l'injection SQL~; (b) l'injection SQL de la page de connexion du portail
des transitaires, version 3.4~; (c) un groupe qui vend des accès à des
réseaux d'entreprise~; (d) la perte possible des données de 4\,000 clients
si le portail est compromis.
\end{boxexo}

\begin{boxexo}[title={\ding{45}\quad Exercice~9.2. Lire un vecteur \niveaumoyen}]
Expliquez en français le vecteur suivant, métrique par métrique~:
\texttt{CVSS:4.0/AV:L/AC:L/AT:N/PR:L/UI:N/VC:H/VI:H/VA:H/SC:N/SI:N/SA:N}.
Cette faille peut-elle servir de point d'entrée depuis Internet~? À quoi
sert-elle typiquement dans une attaque~?
\end{boxexo}

\begin{boxexo}[title={\ding{45}\quad Exercice~9.3. Trois indicateurs \niveaumoyen}]
Une faille a un score CVSS-B de 5,3, un EPSS de 0,92 et figure au catalogue
KEV. Une autre a un score de 9,8, un EPSS de 0,001 et n'y figure pas.
\begin{enumerate}[(a)]
    \item Laquelle corrigez-vous d'abord sur un serveur exposé à Internet~?
    \item Et si les deux étaient sur un serveur interne isolé~?
    \item Que manque-t-il pour décider avec certitude~?
\end{enumerate}
\end{boxexo}

\begin{boxexo}[title={\ding{45}\quad Exercice~9.4. Le contexte de NYANGA \niveauexpert}]
\begin{enumerate}[(a)]
    \item Expliquez pourquoi l'éditeur a coté SC, SI et SA à « nul » pour
          VULN-2024-4471, et pourquoi NYANGA aurait dû les coter
          autrement.
    \item Le 19~décembre 2024, quelle décision SSVC NYANGA aurait-elle dû
          prendre~? Justifiez chaque branche de l'arbre.
    \item Le jour de l'inscription au catalogue, l'EPSS de la faille
          dépassait 0,9 (valeur fictive). Ce chiffre ajoutait-il quelque
          chose à l'inscription~? Discutez.
\end{enumerate}
\end{boxexo}

\bigskip

\begin{boxinfo}
\textbf{Pour aller plus loin.}
\begin{itemize}
    \item La spécification CVSS~4.0~\cite{first2023cvss4}.
    \item Le site de l'EPSS~\cite{firstepss} et le catalogue
          KEV~\cite{cisakev}.
    \item Le guide SSVC de la CISA~\cite{cisa2022ssvc}.
\end{itemize}
\end{boxinfo}

\begin{boxinfo}
\textbf{Transition.} Coter suppose d'avoir trouvé. Comment savoir quelles
vulnérabilités on a, sans attaquer ses propres systèmes~?
Chapitre~\ref{chap:identifier}.
\end{boxinfo}
